\documentclass[10pt,a4paper]{amsart}
\usepackage[margin=19mm]{geometry}
\usepackage{amsmath,amssymb,mathtools}
\usepackage[scr=rsfs,cal=euler]{mathalfa}
\usepackage{xfrac}
\usepackage[unicode,hidelinks,bookmarks=false]{hyperref}
\newcommand{\ie}{i.e.\ }
\newcommand{\cf}{cf.\ }
\newcommand{\resp}{respectively\ }
\newcommand{\Subsection}{\subsection}
\providecommand{\nobreakdash}{\nobreak\hbox{-}}
\setlength{\emergencystretch}{2em}
\multlinegap0pt
\usepackage{amssymb}
\usepackage{mathtools}
\newtheorem{lemma}{Lemma}
\newcommand{\ind}[1]{\mathbf1_{\{#1\}}}
\title{Cyclic Haagerup--Izumi fusion categories\\at every odd order}
\author{Tzu-Chen Huang}
\date{Source: arXiv:2609.15986v1 (CC BY 4.0)}
\begin{document}
\maketitle

\noindent\emph{Source and licence.} Cyclic Haagerup--Izumi fusion categories\\at every odd order. Version arXiv:2609.15986v1 (CC BY 4.0), distributed by arXiv under the Creative Commons Attribution 4.0 International licence, http://creativecommons.org/licenses/by/4.0/, as recorded in the arXiv licence metadata for this exact version and shown on the abstract page https://arxiv.org/abs/2609.15986v1. Full licence notice: \texttt{licenses/arxiv-2609.15986-CC-BY-4.0.txt}.\par\smallskip

\noindent\textit{Editorial scope.} Counted unit: the article's lemma allodd:reflection (source0.tex lines 1716--1720). The article's complete original proof of it is delivered with it (source0.tex lines 1722--1879, one \texttt{proof} environment of 158 lines). No mathematics is written, repaired or re-derived: the statement and the proof are the article's own lines, in the article's own notation, and no retained line is substituted or re-flowed.  Retained uncounted as the source-local context the proof needs: lines 425--432; lines 757--760.  Nothing was omitted from any retained span, so the package carries no omission record.  No retained line contains a citation, so the wrapper carries no bibliography and no bibliography entry.  No retained line contains a figure environment, an image-inclusion command or drawing code, so no image is delivered and the package declares no asset dependency.  Editorial formatting only, in addition: an AMS \texttt{amsart} preamble that restates the article's own declarations and macros used by the retained lines, namely the environments \texttt{lemma} on separate counters, together with the standard packages those lines need.  The retained environments print in this document as Lemma 1 in order of appearance, whereas the article prints the same items with the numbering of its own full document.  The complete native source of the article as distributed in the arXiv e-print for this exact version is bundled unchanged as \texttt{sources/210470/source0.tex}.
\begin{equation}\label{foundation:core}
 \mathsf B_{a,b}=
 \begin{cases}
  W_aW_{b-a}/W_b,&b\ne0,\\
  -1+\tau\ind{a=0},&b=0,
 \end{cases}
 \qquad \mathsf A_{a,b}=\tau^{-1}\mathsf B_{a,b}.
\end{equation}

\begin{equation}\label{tail:uniform-core}
 \mathsf B_{a,b}=\delta^{-1}W_aW_{b-a}W_{-b}
                  +\tau^2\delta^{-1}\ind{a=b=0}.
\end{equation}

\begin{lemma}\label{allodd:reflection}
Suppose the matrix \eqref{foundation:core} satisfies the quadratic equations
and $W_aW_{-a}=\delta$ for $a\ne0$.  Fix $g,h\in G$, put $t=g+h$, and assume $g,h,t\ne0$.  If $C(g,h,k,l)=0$ for every $k,l$, then
$C(-h,-g,k,l)=0$ for every $k,l$.
\end{lemma}

\begin{proof}
We invert the two sides of the cubic block separately. All matrices
have rows and columns indexed by $G$.

\emph{1. Column inverses.} Write $u_p(a)=\mathsf B_{a,p}$ for
$p\ne0$.  The coefficient formula and the quadratic equations give
\begin{equation}\label{allodd:column-inverses}
 u_p(a)=\frac{W_aW_{p-a}}{W_p}=u_p(p-a),\qquad
 \sum_a u_p(a)u_{-p}(x-a)=\tau^2\ind{x=0}.
\end{equation}
For the second identity, multiply $Q(x,p)=0$ by $\tau^2$,
use $\mathsf B_{p,a}=u_{-p}(a-p)=u_{-p}(-a)$, and reindex.
Consequently the circulant with entries $u_p(k-l)$ has inverse
$\tau^{-2}u_{-p}(k-l)$.

\emph{2. The left side.} Define
\[
 L_{g,h}(k,l)=\sum_m u_t(m)\mathsf B_{g,m+k}
                                      \mathsf B_{h,m+l},
 \qquad
 R_{g,h}(k,l)=\mathsf B_{g+l,k}\mathsf B_{h+k,l}.
\]
The assumed cubics say $L_{g,h}=\tau R_{g,h}$.
Let $H_p(k,m)=\mathsf B_{p,k+m}$ and let $D_t$ be the
diagonal matrix with entries $u_t(m)$.  Equation
\eqref{allodd:column-inverses} gives
\[
 H_p^{-1}(k,m)=\tau^{-2}u_p(k+m),\qquad
 L_{g,h}=H_gD_tH_h.
\]
All entries of $D_t$ are nonzero.  Thus $L_{g,h}$ is invertible.
The reciprocal product identities, including the two exceptional samples $m=0,t$,
give
\begin{equation}\label{allodd:diagonal-inverse}
 \frac1{u_t(m)}=\delta^{-1}u_{-t}(-m)
       -\frac{\tau}{\delta}
                    (\ind{m=0}+\ind{m=t}).
\end{equation}
Set
\[
 E(k,l)=u_h(k)u_g(l)+u_h(k+t)u_g(l+t),
 \qquad
 L^{\mathrm r}(k,l)=L_{-h,-g}(-k,-l).
\]
On inserting \eqref{allodd:diagonal-inverse} into
$L_{g,h}^{-1}=H_h^{-1}D_t^{-1}H_g^{-1}$ and replacing
$m$ by $-m$, we obtain
\begin{equation}\label{allodd:left-inverse}
 L_{g,h}^{-1}=\delta^{-1}\tau^{-4}
                         (L^{\mathrm r}-\tau E).
\end{equation}
Here $u_h(k-m)=\mathsf B_{-h,m-k}$, by
\eqref{allodd:column-inverses} and order-three symmetry.

\emph{3. The right side.} We prove, using only the quadratics and
reciprocal products, the corresponding identity
\begin{equation}\label{allodd:right-inverse}
 R^{\mathrm r}=\delta\tau^2 R_{g,h}^{-1}+E,
 \qquad
 R^{\mathrm r}(k,l)=R_{-h,-g}(-k,-l).
\end{equation}
The details below retain the origin corrections in the coefficient formula.
Put
\[
 d_1(k)=W_{h+k}W_{-k},\qquad d_2(l)=W_{g+l}W_{-l},
 \qquad b=u_t(g)=u_t(h)=\frac{W_gW_h}{W_t}\ne0.
\]
Direct substitution of
$\mathsf B_{a,b}=\delta^{-1}W_aW_{b-a}W_{-b}
+\tau^2\delta^{-1}\ind{a=b=0}$ yields
\begin{equation}\label{allodd:right-factor}
 R_{g,h}=\delta^{-2}W_{-t}\,
                   \operatorname{diag}(d_1)M\operatorname{diag}(d_2),
\end{equation}
where
\[
 M(k,l)=u_{-t}(k-l-g)
       +\tau^2(\ind{k=0,l=-g}+\ind{k=-h,l=0}).
\]
The uncorrected matrix $K(k,l)=u_{-t}(k-l-g)$ has inverse
$K^{-1}(k,l)=\tau^{-2}u_t(k-l+g)$.
Writing $U=(e_0,e_{-h})$ and $V=(e_{-g},e_0)$, we have
$M=K+\tau^2 UV^{\mathsf T}$ and
\[
 I_2+\tau^2V^{\mathsf T}K^{-1}U=
       \begin{pmatrix}0&b\\b&0\end{pmatrix}.
\]
This matrix is invertible.  The rank-two inverse formula therefore
proves that $M$, hence $R_{g,h}$, is invertible and that
\begin{equation}\label{allodd:woodbury}
 M^{-1}(k,l)=\tau^{-2}(U_0-U_1-U_2),
\end{equation}
where, for the entry $(k,l)$ in question,
\[
 U_0=u_t(k-l+g),\qquad
 U_1=b^{-1}u_t(k+g)u_t(g-l),\qquad
 U_2=b^{-1}u_t(k+t)u_t(-l).
\]

To check the exceptional entries in
\eqref{allodd:right-inverse}, define
\[
 q_a=W_aW_{-a}=\delta-\tau\ind{a=0},\quad
 \alpha=\frac{q_kq_l}{\delta^2},\quad
 \beta=\frac{q_{g+k}q_{h+l}}{\delta^2},\quad
 Z=\frac{W_{-t}d_2(k)d_1(l)}{\delta^3}.
\]
To see the endpoint corrections explicitly, expand
$R^{\mathrm r}(k,l)=\mathsf B_{-h-l,-k}\mathsf B_{-g-k,-l}$
using \eqref{tail:uniform-core}. The reciprocal products give
\[
\begin{aligned}
 ZR^{\mathrm r}(k,l)
   &=\alpha\beta U_0+\frac{\tau^2}{\delta^2}
       (\ind{k=0,l=-h}+\ind{k=-g,l=0}),\\
 Z\,u_h(k)u_g(l)&=\alpha U_1,\\
 Z\,u_h(k+t)u_g(l+t)&=\beta U_2.
\end{aligned}
\]
In the first line the two corrections come from the two possible
origin entries of $\mathsf B$; they cannot occur together since
$g,h\ne0$. The other two lines follow from $W_tW_{-t}=\delta$
and the definitions of $U_1,U_2$.
Subtracting the two terms of $E$ gives the exact entry identity
\begin{equation}\label{allodd:endpoints}
 Z(R^{\mathrm r}-E)(k,l)
  =\alpha\beta U_0-\alpha U_1-\beta U_2
     +\frac{\tau^2}{\delta^2}
        (\ind{k=0,l=-h}+\ind{k=-g,l=0}).
\end{equation}
If $\alpha\ne1$, then $k=0$ or $l=0$, and $U_1=U_0$.
If $\beta\ne1$, then $k=-g$ or $l=-h$, and $U_2=U_0$.
Both exceptional conditions can hold only at $(k,l)=(0,-h)$
or $(-g,0)$: the assumptions $g,h\ne0$ exclude the other
intersections.  At either point,
$\alpha=\beta=\delta^{-1}$ and $U_0=U_1=U_2=-1$.
Thus the last term of \eqref{allodd:endpoints} cancels
$(\alpha-1)(\beta-1)U_0=-\tau^2/\delta^2$.
In every case the right side of \eqref{allodd:endpoints} is
$U_0-U_1-U_2$.  Equations \eqref{allodd:right-factor} and
\eqref{allodd:woodbury} now give
\[
 Z(R^{\mathrm r}-E)(k,l)
   =\tau^2M^{-1}(k,l)
   =Z\delta\tau^2R_{g,h}^{-1}(k,l).
\]
The factor $Z$ is nonzero, proving
\eqref{allodd:right-inverse} at all entries.

\emph{4. Comparing the inverses.} Now $L_{g,h}=\tau R_{g,h}$ and
\eqref{allodd:left-inverse} imply
\[
 L^{\mathrm r}=\delta\tau^3R_{g,h}^{-1}+\tau E
             =\tau R^{\mathrm r},
\]
where the second equality is \eqref{allodd:right-inverse}.
These are precisely all the cubics with first indices $(-h,-g)$.
\end{proof}

\end{document}
